\documentclass[11pt,a4paper]{article}

\usepackage[margin=1in]{geometry}
\usepackage{amsmath,amssymb,amsthm,mathtools}
\usepackage{enumitem}
\usepackage{hyperref}
\usepackage{booktabs}
\usepackage{array}
\usepackage{longtable}

\newtheorem{theorem}{Theorem}[section]
\newtheorem{proposition}[theorem]{Proposition}
\newtheorem{corollary}[theorem]{Corollary}
\newtheorem{lemma}[theorem]{Lemma}
\newtheorem{definition}[theorem]{Definition}
\newtheorem{remark}[theorem]{Remark}

\newcommand{\Zmod}[1]{\mathbb{Z}/#1\mathbb{Z}}

\title{Galactic Rotation from the Compact Fiber:\\
The Gravitational Limit, $R^{16}$, and a Zero-Parameter Rotation-Curve Model}

\author{A.\ R.\ Wells\\
Dual-Frame Research Group}

\date{\today}

\begin{document}

\maketitle

\begin{abstract}
The coordinate-time mechanism and gravitational limit derived in a companion paper~\cite{Wells2026c} from the postulates of discrete scalar motion predict a characteristic acceleration scale and coherence length at galactic scales.
This paper obtains both quantities from the natural unit system and the inter-regional ratio $R = 156.444$ by dimensional analysis, with no calibrated galactic parameters: $a_0 = a_{\mathrm{nat}} \times N_m / R^{16} = 6.12 \times 10^{-11}$~m/s$^2$ (where $N_m = 4$ is the magnetic sector capacity and $R^{16}$ is the inter-regional exponent identified on independent structural grounds in~\cite{Wells2026c}) and $L_{\mathrm{GL}} = s_{\mathrm{nat}} \times R^{12} = 0.318$~kpc.
The resulting rotation-curve model uses the sqrt-exponential kernel structurally induced by the second-power relation $r = \rho^2$~\cite{Wells2026b} and the covering-degree coupling law from the compact fiber~\cite{Wells2026a} ($d = 2$).
Tested against 171 SPARC galaxies~\cite{Lelli2016}, the model achieves slightly lower median error than a fixed MOND baseline (median $\chi^2$ per point $10.71$ versus $11.25$) while using no calibrated galactic parameters.
The covering degree $d = 2$ is confirmed by sharp empirical selection: departures to $d = 1.8$ or $d = 2.2$ degrade performance by factors of $7$--$100$.
A radial residual analysis shows that the model achieves near-zero mean fractional residuals ($+0.6\%$ to $+0.9\%$) in the mid-to-outer disk, where baryonic-input ambiguities are plausibly reduced relative to the inner disk.
The remaining ${\sim}10\%$ systematic error floor is shared with MOND and traced to galaxy-to-galaxy variation in mass-to-light ratios.
The $R^{16}$ exponent that sets the galactic acceleration scale is the same exponent identified in~\cite{Wells2026c} in the context of the electromagnetic-to-gravitational hierarchy ratio, providing a cross-scale consistency check between atomic and galactic physics.

This paper does not claim that the SPARC benchmark proves the compact-fiber framework; it claims that the framework, as developed in~\cite{Wells2026a,Wells2026b,Wells2026c}, yields a specific galactic rotation model with no calibrated galactic parameters, and that this model performs competitively with standard MOND on the SPARC sample.
The coordinate-time metric passed Solar System tests in~\cite{Wells2026c}; the additive nonlocal galactic term introduced here is argued to be independently safe: explicit numerical evaluation shows that the nonlocal kernel evaluates $\tilde{g}$ at galactic scales, producing a constant extra acceleration across the Solar System that is absorbed by free fall, with a residual galactic tidal contribution of ${\sim}10^{-16}$ of Solar gravity (Appendix~\ref{app:solar-numerical}).
\end{abstract}

\tableofcontents

%===================================================================
\section{Introduction}
\label{sec:intro}
%===================================================================

\subsection{Aim}

A companion paper~\cite{Wells2026c} derived the coordinate-time mechanism for gravitation from the postulates of discrete scalar motion, demonstrated competency via Solar System tests ($\beta = \gamma = 1$, four classical tests, the equivalence principle), and identified two open items at status~\textbf{O}: the gravitational limit and the $R^{16}$ acceleration scale.
This paper develops both into an explicit galactic rotation model and evaluates their empirical viability on the SPARC sample.
It obtains the galactic acceleration scale and coherence length from natural units and the compact fiber structure, constructs an explicit rotation-curve model, and tests it against the SPARC sample~\cite{Lelli2016}.
All derivations rest on the postulates and structures established in~\cite{Wells2026a,Wells2026b,Wells2026c}; we import only the items listed in~\S\ref{sec:dependency}.

The galactic model is a conditional consequence of Papers~I--III plus one auxiliary assumption: exponential smoothing in the time-region radial coordinate~$\rho$.
It is not a uniquely forced derivation from first principles.
The auxiliary assumption is structurally motivated (\S\ref{sec:kernel}) and empirically confirmed as optimal (\S\ref{sec:kernel-scan}), but alternative smoothing forms are not excluded by the postulates alone.
This conditionality is reflected in the status labels throughout.

\subsection{What this paper does and does not do}
\label{sec:scope}

This paper:
\begin{itemize}
\item derives $a_0$ and $L_{\mathrm{GL}}$ from the natural unit system and the inter-regional ratio~$R$,
\item constructs an explicit rotation-curve model with zero calibrated galactic parameters,
\item tests against 171 SPARC galaxies with a fixed preprocessing pipeline and explicit comparison protocol,
\item analyzes failure modes and identifies the next structural target,
\item reports sensitivity to implementation choices (kernel power, curvature modulation, mass-to-light ratio).
\end{itemize}

This paper does not:
\begin{itemize}
\item derive $G$ from natural units (still~\textbf{O}),
\item treat cosmological implications (CMB, BAO, large-scale structure),
\item claim to replace $\Lambda$CDM,
\item fit per-galaxy parameters,
\item claim that the SPARC benchmark proves the compact-fiber framework.
\end{itemize}

\noindent In this paper, ``zero calibrated galactic parameters'' means that no parameters in the RST law are tuned on a per-galaxy basis or fitted to the SPARC sample; standard SPARC baryonic inputs ($\Upsilon_{\mathrm{disk}}$, $\Upsilon_{\mathrm{bulge}}$) are held fixed across all models as external observational inputs rather than calibrated theory parameters.

\subsection{Dependency on companion papers}
\label{sec:dependency}

From~\cite{Wells2026a} (compact fiber):
\begin{itemize}
\item Fiber capacities $N_m = 4$ (magnetic), $N_e = 8$ (electric).
\item Covering degree $d = 2$.
\item Transport metric $1/N$ per step (Theorem~12.1 of~\cite{Wells2026a}).
\end{itemize}

From~\cite{Wells2026b} (atomic structure):
\begin{itemize}
\item Second-power relation $r = \rho^2$ (Proposition~7.1 of~\cite{Wells2026b}).
\item The structural constant $6 = \mathrm{period}_{\mathrm{ref}} + e_{\mathrm{half}}$ (Proposition~6.4 of~\cite{Wells2026b}, the constant $42 = 6 \times 7$).
\item Natural units $s_{\mathrm{nat}} = 4.559 \times 10^{-8}$~m, $t_{\mathrm{nat}} = 1.521 \times 10^{-16}$~s.
\item Inter-regional ratio $R = 156.444$.
\end{itemize}

From~\cite{Wells2026c} (quantum and gravitational structure):
\begin{itemize}
\item Coordinate-time mechanism ($3v^2/c^2$ per clock time).
\item Gravitational limit concept (\S14.9 of~\cite{Wells2026c}).
\item Reciprocal partition and effective metric (\S14.3 of~\cite{Wells2026c}).
\item $R^{16}$ identified in hierarchy-ratio discussion (\S14.10 of~\cite{Wells2026c}).
\item PPN parameters $\beta = \gamma = 1$; Solar System tests passed.
\end{itemize}

\subsection{Relation to existing approaches}
\label{sec:existing}

MOND~\cite{Milgrom1983} provides an empirical interpolation function with one calibrated parameter ($a_0 = 1.2 \times 10^{-10}$~m/s$^2$) but no structural derivation of either $a_0$ or the interpolation form.
Dark matter halos (NFW, Burkert) require 2--3 per-galaxy parameters and do not predict the baryonic Tully-Fisher relation from first principles.
The radial acceleration relation (RAR)~\cite{McGaugh2016} is an empirical phenomenological regularity.
The present paper offers a structural derivation from discrete scalar motion postulates, with zero calibrated parameters and specific predictions about where the model should fail.

\subsection{Status discipline}
\label{sec:status}

Every model ingredient is tagged with one of:
\begin{description}[style=nextline,leftmargin=2.5cm]
\item[\textbf{T} --- Theorem-derived] Complete chain from imported structure. No gaps.
\item[\textbf{S} --- Structurally inferred] Follows under explicitly stated assumptions.
\item[\textbf{S\!f} --- Implementation-fixed] Structurally motivated, frozen before comparison, shown to be non-critical by sensitivity analysis.
\item[\textbf{E} --- Empirically tested] Quantitative comparison with data.
\item[\textbf{O} --- Open] Identified gap.
\end{description}

The additional label~\textbf{S\!f} distinguishes implementation choices that are frozen before comparison from quantities that are derived from the postulates.

%===================================================================
\section{The Gravitational Limit}
\label{sec:grav-limit}
%===================================================================

\subsection{Derivation from the postulates}

The scalar progression is universal outward motion at unit speed~$c$ (P5).
Gravitation is three-dimensional scalar motion inward with speed $v^2 = GM/r$, decreasing with distance.
At some distance $r_{\mathrm{GL}}$ from a mass concentration, the two speeds are equal; beyond $r_{\mathrm{GL}}$, the net scalar motion is outward.
This is not a modification of gravity.
The progression exists at all scales but is dynamically negligible where gravitation dominates (the Solar System) and becomes significant where gravitation is weak (galactic outskirts).
\textit{Status:~\textbf{S}} (structural consequence of P5 and the gravitational speed profile).

\subsection{What the gravitational limit replaces}

In the standard framework, the discrepancy between observed rotation curves and Newtonian predictions is attributed to a dark matter halo providing extra mass.
In MOND, it is attributed to a modified gravitational law below the acceleration scale~$a_0$.
In the RST, it is attributed to neither: the ever-present scalar progression becomes dynamically significant at the gravitational limit, producing the observed flat rotation curves without extra mass or modified gravity.

\subsection{Solar System safety}

The nonlocal correction at Solar System radii is treated in~\S\ref{sec:solar}.
In brief: the nonlocal kernel evaluates $\tilde{g}$ at galactic scales (not Solar System scales), producing a constant extra acceleration $g_{\mathrm{extra}} \approx 1.1 \times 10^{-10}$~m/s$^2$ across the entire Solar System.
By the equivalence principle, this constant acceleration is undetectable in the freely falling Solar System frame.
The residual galactic tidal contribution is ${\sim}10^{-16}$ of Solar gravity.
The detailed argument is given in~\S\ref{sec:solar}.

%===================================================================
\section{Derivation of the Acceleration Scale}
\label{sec:accel-scale}
%===================================================================

\subsection{The natural acceleration}

The natural acceleration is
\begin{equation}
\label{eq:a-nat}
a_{\mathrm{nat}} = \frac{s_{\mathrm{nat}}}{t_{\mathrm{nat}}^2} = \frac{4.559 \times 10^{-8}}{(1.521 \times 10^{-16})^2} = 1.97 \times 10^{24} \;\text{m/s}^2.
\end{equation}
This is the only acceleration scale available from the postulates before inter-regional conversion.

\subsection{Why $R^{16}$}
\label{sec:R16}

Mass has RST dimensions $t^3/s^3$.
Converting $a_{\mathrm{nat}}$ to conventional gravitational scales requires the inter-regional ratio~$R$ for each dimensional crossing.
In~\cite{Wells2026c} (\S14.10), $R^{16} \approx 1.29 \times 10^{35}$ was identified on independent structural grounds as the power arising in the dimensional analysis connecting natural units to gravitational scales, in the context of the hierarchy ratio $\alpha/\alpha_G \approx 10^{36}$.

The present paper examines one specific downstream consequence: the galactic acceleration scale $a_{\mathrm{nat}}/R^{16}$.
The SPARC comparison tests whether this consequence is empirically viable, not whether $R^{16}$ is correct.
The structural argument for the exponent~16 is given in~\cite{Wells2026c}; here it is imported, not re-derived.

\textit{Status:~\textbf{S}} (imported from~\cite{Wells2026c} under stated interpretation of the dimensional chain).

\subsection{The magnetic-sector factor}
\label{sec:Nm}

Gravitation is three-dimensional scalar motion, which occupies the magnetic sector of the compact fiber~\cite{Wells2026a}.
The magnetic sector capacity $N_m = 4$ is therefore the structural factor entering the gravitational acceleration scale:
\begin{equation}
\label{eq:a0}
a_0 = \frac{a_{\mathrm{nat}} \times N_m}{R^{16}} = \frac{1.97 \times 10^{24} \times 4}{1.29 \times 10^{35}} = 6.12 \times 10^{-11} \;\text{m/s}^2.
\end{equation}

The choice of $N_m$ follows from the sector assignment: gravitation is magnetic-sector motion, so the gravitational coupling involves the magnetic capacity.
This is the same $N_m$ that determines the cross-sector screening slopes in~\cite{Wells2026b} (ionization energy, electron affinity, polarizability: all have slope $1/N_m = 1/4$).
The factor is not selected by empirical performance; it is determined by the structural role of the magnetic sector in gravitation.

\textit{Status:~\textbf{S}} (structural argument for $N_m$ as the gravitational sector factor).

\subsection{Contextual comparison with MOND}
\label{sec:mond-context}

The conventional MOND acceleration scale is $a_0^{\mathrm{MOND}} = 1.2 \times 10^{-10}$~m/s$^2$~\cite{Milgrom1983}, an empirically calibrated value.
The RST-derived scale $a_0 = 6.12 \times 10^{-11}$~m/s$^2$ is approximately half the MOND value.

It is noted that applying the electric sector capacity $N_e = 8$ in place of $N_m$ gives $a_{\mathrm{nat}} \times N_e / R^{16} = 1.22 \times 10^{-10}$~m/s$^2$, within 2\% of the MOND value.
This proximity is recorded as a contextual observation: the conventional MOND scale may correspond to the electric-sector reading of the same dimensional chain.
It plays no role in the model specification.

%===================================================================
\section{Derivation of the Coherence Length}
\label{sec:coherence}
%===================================================================

\subsection{$L_{\mathrm{GL}}$ from natural units and $R^{12}$}

\begin{equation}
\label{eq:LGL}
L_{\mathrm{GL}} = s_{\mathrm{nat}} \times R^{12} = 4.559 \times 10^{-8} \times (156.444)^{12} = 0.318 \;\text{kpc}.
\end{equation}

The power~12 decomposes as $2 \times 6$, where the factor~2 is $R^2$ (one inter-regional crossing) and the factor $6 = \mathrm{period}_{\mathrm{ref}} + e_{\mathrm{half}}$ is the same structural constant that enters the $n_{\mathrm{eff}}$ law in~\cite{Wells2026b} (Proposition~6.4, the constant $42 = 6 \times 7$).
This identification is structurally consistent but has not been derived from a closed dimensional chain.

\textit{Status:~\textbf{S}} (dimensional analysis plus structural identification of the factor~6).

\subsection{Physical interpretation}

$L_{\mathrm{GL}}$ is the coherence scale of the gravitational-limit transition: the radial distance over which the baryonic field is nonlocally averaged to evaluate the progression-gravitation balance.
It is not the gravitational limit distance itself (which depends on the enclosed mass profile and varies from galaxy to galaxy).
It governs the sharpness of the transition from Newtonian-dominated to progression-dominated dynamics.

%===================================================================
\section{The RST Rotation Model}
\label{sec:model}
%===================================================================

\subsection{Kernel shape: structurally induced by the second-power relation}
\label{sec:kernel}

The second-power relation $r = \rho^2$~(\cite{Wells2026b}, Proposition~7.1) maps the time-region radial coordinate~$\rho$ to the extension-space coordinate~$r$.
An exponential smoothing kernel in the time region,
\begin{equation}
K(\rho, \rho') = \exp\!\left(-\frac{|\rho - \rho'|}{L_\rho}\right),
\end{equation}
maps under $r = \rho^2$ to a sqrt-exponential kernel in extension space:
\begin{equation}
\label{eq:kernel}
K(r, r') = \exp\!\left(-\frac{|\sqrt{r} - \sqrt{r'}|}{L_{\mathrm{GL}}}\right).
\end{equation}

This is the natural induced kernel under the assumption of exponential smoothing in the time region.
It is not the unique possible kernel: other time-region smoothing forms would induce different extension-space kernels.
The assumption of exponential smoothing is motivated by the isotonic structure of the time-region potential~(\cite{Wells2026b}, \S7.3) but is not derived from it as a theorem.

The kernel power scan (\S\ref{sec:kernel-scan}) confirms that the structurally induced power $p = 0.5$ lies at the empirical optimum: performance degrades monotonically for $p > 0.6$, and the linear kernel ($p = 1.0$) performs worse than MOND.

\textit{Status:~\textbf{S}} (structurally induced under a stated assumption; empirically confirmed as optimal).

\subsection{Coupling law from the covering degree}
\label{sec:coupling}

The covering degree $d = 2$~(\cite{Wells2026a}, Theorem~7.4) determines the coupling exponent in the effective acceleration law:
\begin{equation}
\label{eq:coupling}
g_{\mathrm{extra}}(r) = \bigl[a_0 \cdot \tilde{g}(r)\bigr]^{1/d} = \sqrt{a_0 \cdot \tilde{g}(r)}.
\end{equation}

This is the unique exponent consistent with the degree-2 covering projection.

\textit{Status:~\textbf{T}} (from~\cite{Wells2026a}).
The empirical confirmation (\S\ref{sec:d-scan}) shows $d = 2.0$ sharply selected; departures to $d \neq 2$ degrade performance by factors of 7--100.

\subsection{Full model specification (pure structural form)}
\label{sec:model-spec}

The pure structural RST rotation model is:
\begin{align}
g_{\mathrm{bar}}(r) &= \frac{v_{\mathrm{bar}}^2(r)}{r}, \label{eq:gbar} \\[4pt]
\tilde{g}(r) &= \frac{\sum_j g_{\mathrm{bar}}(r_j)\,K(r, r_j)}{\sum_j K(r, r_j)}, \label{eq:gtilde} \\[4pt]
g_{\mathrm{total}}(r) &= g_{\mathrm{bar}}(r) + \sqrt{a_0 \cdot \tilde{g}(r)}, \label{eq:gtotal} \\[4pt]
v_{\mathrm{pred}}(r) &= \sqrt{g_{\mathrm{total}}(r) \cdot r}, \label{eq:vpred}
\end{align}
where $K(r,r')$ is given by Eq.~\eqref{eq:kernel} and $v_{\mathrm{bar}}$ is the Newtonian baryonic velocity from SPARC components (gas, disk, bulge).

Parameters derived from~\cite{Wells2026a,Wells2026b,Wells2026c}:
\begin{center}
\begin{tabular}{llll}
\toprule
\textbf{Quantity} & \textbf{Value} & \textbf{Source} & \textbf{Status} \\
\midrule
$a_0$ & $6.12 \times 10^{-11}$ m/s$^2$ & $a_{\mathrm{nat}} \times N_m / R^{16}$ & \textbf{S} \\
$L_{\mathrm{GL}}$ & $0.318$ kpc & $s_{\mathrm{nat}} \times R^{12}$ & \textbf{S} \\
$d$ & $2$ & Covering degree & \textbf{T} \\
$p$ (kernel power) & $0.5$ & From $r = \rho^2$ & \textbf{S} \\
\bottomrule
\end{tabular}
\end{center}

Calibrated parameters: none.
Implementation constants: none.

This is the headline model.
All ingredients are traced to~\cite{Wells2026a,Wells2026b,Wells2026c} with explicit status labels.

\subsection{Ingredient provenance}
\label{sec:provenance}

\begin{center}
\begin{tabular}{llll}
\toprule
\textbf{Ingredient} & \textbf{Source} & \textbf{Status} & \textbf{Note} \\
\midrule
$a_0$ & \cite{Wells2026c} & \textbf{S} & $a_{\mathrm{nat}} \times N_m / R^{16}$ \\
$L_{\mathrm{GL}}$ & \cite{Wells2026b} & \textbf{S} & $s_{\mathrm{nat}} \times R^{12}$ \\
$d = 2$ & \cite{Wells2026a} & \textbf{T} & Covering degree \\
Kernel power $p = 0.5$ & \cite{Wells2026b} & \textbf{S} & From $r = \rho^2$ \\
Kernel form (exponential in $\rho$) & assumed & \textbf{S} & Isotonic potential motivation \\
$\Upsilon_{\mathrm{disk}} = 0.5$ & SPARC standard & \textbf{S\!f} & 3.6~$\mu$m M/L \\
$\Upsilon_{\mathrm{bulge}} = 0.7$ & SPARC standard & \textbf{S\!f} & 3.6~$\mu$m M/L \\
Preprocessing & SPARC standard & \textbf{S\!f} & No galaxy exclusions \\
\bottomrule
\end{tabular}
\end{center}

\subsection{Optional refinement: curvature-modulated coherence length}
\label{sec:curvature}

The transport metric $1/N$ per step~(\cite{Wells2026a}, Theorem~12.1) motivates a radially varying coherence length that shortens where the baryonic profile has sharp features and lengthens where it is smooth:
\begin{equation}
\label{eq:Leff}
L_{\mathrm{eff}}(r) = L_{\mathrm{GL}} \times \mathrm{clip}\!\left(1 + \eta \cdot \tanh\!\left(\hat{C}(r)\right),\; 1 - \eta,\; 1 + \eta\right),
\end{equation}
where $\hat{C}(r)$ is the normalized curvature of $\log g_{\mathrm{bar}}$ along radius.
This is motivated by the fiber transport metric but is not derived from it as a theorem.
The modulation depth $\eta$ is an implementation constant.

The sensitivity analysis (Table~\ref{tab:eta}) shows the headline result is insensitive to~$\eta$:

\begin{table}[h]
\centering
\caption{Curvature modulation sensitivity. The pure structural model ($\eta = 0$) is the headline; curvature modulation provides a marginal refinement.}
\label{tab:eta}
\begin{tabular}{rrl}
\toprule
$\eta$ & Median $\chi^2$/pt & vs.\ MOND \\
\midrule
0.00 & 10.71 & 0.951 \\
0.20 & 10.59 & 0.941 \\
0.45 & 10.43 & 0.927 \\
0.55 & 10.36 & 0.921 \\
0.70 & 10.66 & 0.948 \\
\bottomrule
\end{tabular}
\end{table}

The entire range $\eta = 0$ to $\eta = 0.7$ performs competitively with MOND.
The improvement from $\eta = 0$ to $\eta = 0.45$ is 2.6\%.
The curvature modulation helps at the margin but is not responsible for the model's competitiveness.

\textit{Status:~\textbf{S\!f}} (structurally motivated, implementation-fixed, non-critical by sensitivity scan).

\subsection{Comparison models}
\label{sec:comparison-models}

\begin{center}
\begin{tabular}{llll}
\toprule
\textbf{Model} & \textbf{Kernel} & \textbf{Parameters} & \textbf{Role} \\
\midrule
MOND/RAR & algebraic & 1 (empirical) & benchmark \\
DFT-B & triangular (empirical) & 2 (calibrated) & ablation comparator \\
RST & $\sqrt{\phantom{x}}$-exp (structurally induced) & 0 & this paper \\
\bottomrule
\end{tabular}
\end{center}

MOND uses the exponential-style RAR mapping $g_{\mathrm{total}} = g_{\mathrm{bar}} / (1 - e^{-\sqrt{g_{\mathrm{bar}}/a_0^{\mathrm{MOND}}}})$ with $a_0^{\mathrm{MOND}} = 1.2 \times 10^{-10}$~m/s$^2$~\cite{Milgrom1983,McGaugh2016}.
DFT-B uses the same coupling law as RST ($g + \sqrt{a_c \cdot \tilde{g}}$) but with a compact triangular kernel $K = \max(0, 1 - |r - r'|/L_c)$ and two calibrated parameters $(a_c, L_c)$.
DFT-B is an intermediate nonlocal benchmark, constructed independently of the RST theoretical framework, that shares the coupling functional form but uses an empirically chosen kernel and calibrated parameters.
It is included to isolate the contribution of the RST kernel derivation and parameter-free specification: any performance difference between DFT-B and RST is attributable to the kernel shape and parameter source, not to the coupling law.

%===================================================================
\section{Empirical Results}
\label{sec:results}
%===================================================================

\subsection{Dataset and preprocessing}

The dataset is the SPARC Rotmod\_LTG catalog~\cite{Lelli2016}: observed rotation velocities, baryonic velocity components (gas, disk, bulge), and measurement uncertainties for late-type galaxies.
From the full catalog, 171 galaxies with at least 5 data points and positive velocity uncertainties are retained.
No galaxies are excluded by morphological, quality, or performance criteria.
Mass-to-light ratios are fixed at the SPARC standard values for 3.6~$\mu$m photometry: $\Upsilon_{\mathrm{disk}} = 0.5\,M_\odot/L_\odot$, $\Upsilon_{\mathrm{bulge}} = 0.7\,M_\odot/L_\odot$.
No preprocessing is applied beyond sorting by radius.

The comparison metric is $\chi^2$ per point:
\begin{equation}
\chi^2/\mathrm{pt} = \frac{1}{N} \sum_{i=1}^{N} \left(\frac{v_{\mathrm{obs},i} - v_{\mathrm{model},i}}{\sigma_{v,i}}\right)^2,
\end{equation}
where $\sigma_{v,i}$ is the reported velocity uncertainty.

\subsection{Global performance}
\label{sec:global}

\begin{table}[h]
\centering
\caption{Global performance on 171 SPARC galaxies. The RST model uses zero calibrated galactic parameters.}
\label{tab:global}
\begin{tabular}{lcrrrrrr}
\toprule
\textbf{Model} & \textbf{Params} & \textbf{Median} & \textbf{p25} & \textbf{p75} & \textbf{p90} & \textbf{p95} & \textbf{Wins} \\
\midrule
MOND & 1 & 11.25 & 3.50 & 35.35 & 90.98 & 155.4 & --- \\
DFT-B & 2 & 10.13 & --- & 41.06 & 108.4 & 150.9 & 85/171 \\
RST & 0 & 10.71 & 3.48 & 41.74 & 96.70 & 149.3 & 87/171 \\
\bottomrule
\end{tabular}
\end{table}

The RST model achieves slightly lower median error than MOND ($10.71$ vs.\ $11.25$, ratio $0.951$) while using no calibrated galactic parameters.
The galaxy-level win count ($87/171$, $51\%$) is consistent with a population-level statistical tie (sign test $p = 0.65$), with the structural advantage being the absence of calibrated parameters.

\subsection{Distribution analysis}
\label{sec:distribution}

\begin{table}[h]
\centering
\caption{Distribution of $\chi^2$/pt across 171 galaxies.}
\label{tab:dist}
\begin{tabular}{lrr}
\toprule
\textbf{Percentile} & \textbf{RST} & \textbf{MOND} \\
\midrule
p10 & 1.49 & 1.99 \\
p25 & 3.48 & 3.50 \\
p50 & 10.71 & 11.25 \\
p75 & 41.74 & 35.35 \\
p90 & 96.70 & 90.98 \\
p95 & 149.3 & 155.4 \\
\midrule
$\chi^2 < 2$ (excellent) & 25 & 18 \\
$\chi^2 < 5$ (good) & 59 & 56 \\
\bottomrule
\end{tabular}
\end{table}

RST produces more excellent fits ($\chi^2 < 2$: 25 vs.\ 18) and more good fits ($\chi^2 < 5$: 59 vs.\ 56) than MOND.
Both models show comparable tail behavior (p90 and p95 within $\sim\!7\%$).
The median log-ratio $\log_{10}(\chi^2_{\mathrm{RST}} / \chi^2_{\mathrm{MOND}})$ is $-0.022$, indicating a slight systematic preference for RST.

The models are globally competitive, but performance differs across galaxy populations: tied overall, systematically better in low-mass dwarfs and in mid-to-outer disk residuals (\S\ref{sec:radial}), weaker in massive disk-dominated spirals (\S\ref{sec:velocity}).

\subsection{Covering degree confirmation}
\label{sec:d-scan}

\begin{table}[h]
\centering
\caption{Covering-degree scan from the canonical model state. All parameters except $d$ are held at their derived values.}
\label{tab:d-scan}
\begin{tabular}{rrr}
\toprule
$d$ & Median $\chi^2$/pt & Ratio to MOND \\
\midrule
1.5 & 85.6 & 7.6$\times$ \\
1.8 & 74.3 & 6.6$\times$ \\
\textbf{2.0} & \textbf{10.71} & \textbf{0.95}$\times$ \\
2.2 & 1129 & 100$\times$ \\
2.5 & 31\,660 & 2814$\times$ \\
3.0 & 985\,442 & 87\,573$\times$ \\
\bottomrule
\end{tabular}
\end{table}

The covering degree $d = 2$ is sharply selected by the SPARC benchmark: performance degrades by over an order of magnitude for modest departures from $d = 2$, and by several orders of magnitude for larger departures.
Within the model family defined by the RST coupling law, $d = 2$ is sharply and uniquely selected by the SPARC benchmark.
This is supportive of the compact-fiber structure~\cite{Wells2026a} insofar as the derived covering degree is sharply selected within the proposed model family, though it does not by itself constitute an independent confirmation of the underlying ontology.

\subsection{Kernel power confirmation}
\label{sec:kernel-scan}

\begin{table}[h]
\centering
\caption{Kernel power scan from the canonical model state. All parameters except $p$ are held at their derived values.}
\label{tab:kernel-scan}
\begin{tabular}{rrrl}
\toprule
$p$ & Median $\chi^2$/pt & Ratio to MOND & Note \\
\midrule
0.30 & 11.15 & 0.991 & \\
0.40 & 10.74 & 0.955 & \\
\textbf{0.50} & \textbf{10.71} & \textbf{0.951} & $r = \rho^2$ (structural) \\
0.60 & 11.16 & 0.992 & \\
0.70 & 11.33 & 1.007 & loses to MOND \\
0.80 & 11.59 & 1.030 & \\
1.00 & 11.80 & 1.048 & linear kernel \\
\bottomrule
\end{tabular}
\end{table}

The structurally derived power $p = 0.5$ lies at the empirical optimum.
Powers above $0.6$ lose to MOND.
The linear kernel ($p = 1.0$) is clearly worse.
The second-power relation~\cite{Wells2026b}, together with exponential smoothing in $\rho$-space, selects the empirically preferred kernel within the tested family.

%===================================================================
\section{Performance Ceiling and Baryonic Input Limitation}
\label{sec:ceiling}
%===================================================================

\subsection{The systematic error floor}

Both RST and MOND produce $\chi^2/\mathrm{pt} \approx 10$, roughly $10\times$ the measurement noise floor ($\chi^2/\mathrm{pt} = 1$ for a perfect model with correctly estimated uncertainties).
The median fractional measurement uncertainty is $\sigma_v / v_{\mathrm{obs}} \approx 5\%$.
A ${\sim}15\%$ systematic model error at $5\%$ measurement precision produces $\chi^2/\mathrm{pt} \approx (0.15/0.05)^2 = 9$, consistent with the observed values.
The systematic residual is shared by both models and is not a feature of the RST framework specifically.

\subsection{Mass-to-light ratio sensitivity}
\label{sec:ml-sensitivity}

\begin{table}[h]
\centering
\caption{Sensitivity of the RST pure structural model to the global disk mass-to-light ratio.}
\label{tab:ml}
\begin{tabular}{rr}
\toprule
$\Upsilon_{\mathrm{disk}}$ & Median $\chi^2$/pt \\
\midrule
0.3 & 14.6 \\
0.4 & 12.2 \\
\textbf{0.5} & \textbf{10.7} \\
0.6 & 11.2 \\
0.7 & 14.8 \\
\bottomrule
\end{tabular}
\end{table}

Varying $\Upsilon_{\mathrm{disk}}$ by $\pm 0.1$ (a 20\% change in disk mass) changes the median $\chi^2$ by 15--40\%.
The baryonic mass model appears to be the leading identified source of systematic error, rather than the gravitational theory.
MOND is subject to the same limitation: its $a_0$ was calibrated assuming $\Upsilon_{\mathrm{disk}} = 0.5$.
Neither model can exceed the performance ceiling set by galaxy-to-galaxy M/L scatter without per-galaxy baryonic fitting.

\subsection{Radial residual profile}
\label{sec:radial}

\begin{table}[h]
\centering
\caption{Mean fractional residual $(v_{\mathrm{model}} - v_{\mathrm{obs}})/v_{\mathrm{obs}}$ by normalized radius, averaged over all per-point measurements across 171 galaxies.}
\label{tab:radial}
\begin{tabular}{rrrrrr}
\toprule
$r/r_{\max}$ & RST mean & RST rms & MOND mean & MOND rms \\
\midrule
$0.0$--$0.2$ & $+12.6\%$ & $40.5\%$ & $+10.6\%$ & $46.5\%$ \\
$0.2$--$0.4$ & $+3.3\%$ & $27.2\%$ & $+4.9\%$ & $29.7\%$ \\
$0.4$--$0.6$ & $+0.9\%$ & $22.3\%$ & $+4.8\%$ & $24.6\%$ \\
$0.6$--$0.8$ & $+0.6\%$ & $20.7\%$ & $+5.8\%$ & $22.7\%$ \\
$0.8$--$1.0$ & $+3.5\%$ & $22.4\%$ & $+8.9\%$ & $24.3\%$ \\
\bottomrule
\end{tabular}
\end{table}

RST achieves near-zero mean residuals in the mid-to-outer disk ($0.4$--$0.8\,r_{\max}$: point-pooled mean $+0.6\%$ to $+0.9\%$), the region where baryonic-input ambiguities are plausibly reduced relative to the inner disk.
MOND systematically overpredicts by $5$--$9\%$ in the same region.
A galaxy-level robustness check (Appendix~\ref{app:robustness}), which computes per-galaxy mean residuals and then takes the median across galaxies, confirms the same pattern: RST medians of $-2.7\%$ to $+0.1\%$ in the $0.2$--$0.8\,r_{\max}$ range, versus MOND medians of $+0.3\%$ to $+6.1\%$.
Both models overpredict in the inner region ($+10$--$13\%$), which is dominated by baryonic input sensitivity ($\Upsilon_{\mathrm{disk}}$, bulge model).

In the mid-to-outer disk, the remaining room for improvement is limited by baryonic input uncertainty; a systematic residual reduction of this scale is therefore nontrivial relative to the available headroom.

%===================================================================
\section{Failure Mode Analysis}
\label{sec:failures}
%===================================================================

\subsection{Velocity-dependent performance}
\label{sec:velocity}

\begin{table}[h]
\centering
\caption{RST performance by galaxy velocity class.}
\label{tab:velocity}
\begin{tabular}{rrrr}
\toprule
$V_{\mathrm{flat}}$ (km/s) & $n$ & RST win rate & Median $\log_{10}$ ratio \\
\midrule
$0$--$50$ & 23 & 100\% & $-0.181$ \\
$50$--$100$ & 63 & 57\% & $-0.032$ \\
$100$--$150$ & 31 & 29\% & $+0.071$ \\
$150$--$250$ & 43 & 35\% & $+0.109$ \\
$250$--$400$ & 11 & 55\% & $-0.008$ \\
\bottomrule
\end{tabular}
\end{table}

RST dominates in low-mass, gas-rich dwarfs ($V_{\mathrm{flat}} < 50$~km/s: 100\% win rate) and is competitive in giant spirals ($V_{\mathrm{flat}} > 250$~km/s).
The deficit is concentrated in disk-dominated massive spirals ($V_{\mathrm{flat}} = 100$--$250$~km/s).

\subsection{Structural interpretation}

Gas-dominated dwarfs have smooth, extended baryonic profiles.
The sqrt-exponential kernel's heavy tails correctly capture the gradual gravitational-limit transition in these systems.
Disk-dominated massive spirals have sharp radial features (bulge--disk transition, disk truncation) that the global coherence length $L_{\mathrm{GL}}$ cannot track.
A mass-dependent gravitational-limit distance (larger for more massive systems) would naturally improve this regime.
This is identified as the next structural target for the program.

\subsection{What would falsify the model}
\label{sec:falsification}

The model makes specific falsifiable predictions:
\begin{enumerate}
\item Any clean-interior dwarf galaxy ($V_{\mathrm{flat}} < 50$~km/s, no morphological complications) where RST fails with $\chi^2 > 15$ would indicate a structural failure, not a baryonic input issue.
\item A systematic trend in residuals uncorrelated with galaxy properties would indicate a missing structural mechanism.
\item A covering degree $d \neq 2$ required for good fits would falsify the compact fiber.
\item If a future larger sample shifted the kernel-power optimum away from $p = 0.5$, the second-power relation would be in tension.
\end{enumerate}

%===================================================================
\section{Cross-Scale Consistency: $R^{16}$ at Atomic and Galactic Scales}
\label{sec:cross-scale}
%===================================================================

\subsection{$R^{16}$ in the companion paper}

The companion paper~\cite{Wells2026c} (\S14.10) identified $R^{16} \approx 1.29 \times 10^{35}$ in the dimensional analysis connecting natural units to gravitational scales, in the context of the hierarchy ratio $\alpha/\alpha_G \approx 10^{36}$.
That identification was made on structural grounds independent of galactic data.

\subsection{$R^{16}$ in this paper}

The galactic acceleration scale $a_0 = a_{\mathrm{nat}} \times N_m / R^{16}$ uses the same exponent.
This is not a validation of $R^{16}$ from the galactic data---the exponent was fixed before the SPARC comparison.
It is a cross-scale consistency check: the same structural quantity that appears in the hierarchy discussion also produces a viable galactic acceleration scale.

\subsection{Structural decomposition}

$R^{16} = (R^2)^8$: $R^2$ per dimensional crossing, iterated $N_e = 8$ times.
$R^{12} = (R^2)^6$: $R^2$ per crossing, iterated $6 = \mathrm{period}_{\mathrm{ref}} + e_{\mathrm{half}}$ times.
These decompositions are structurally consistent but not derived from a closed chain.
The exact derivation of why the hierarchy involves $R^{16}$ rather than a nearby power requires deriving $G$ from natural units (status~\textbf{O}).

\subsection{What closing the $G$ derivation would add}

If $G$ were derived from natural units as $G = f(s_{\mathrm{nat}}, t_{\mathrm{nat}}, R)$, the hierarchy dissolution~(\cite{Wells2026c}, \S14.10) would close, and the galactic $a_0$ would follow as a consequence of the $G$ derivation rather than as an independent dimensional inference.
\textit{Status:~\textbf{O}.}

%===================================================================
\section{Solar System Safety}
\label{sec:solar}
%===================================================================

The RST model adds a correction $\sqrt{a_0 \cdot \tilde{g}(r)}$ to the Newtonian baryonic acceleration.
This section shows that this correction does not produce detectable effects on planetary orbits, through two independent mechanisms.

\subsection{Mechanism 1: Nonlocal averaging evaluates the galactic field, not the Solar field}

The nonlocal kernel (Eq.~\ref{eq:kernel}) averages the baryonic field over a coherence scale $L_{\mathrm{GL}} = 0.318$~kpc $= 65\,600$~AU.
This is far larger than the Sun's gravitational dominance sphere: the Sun's gravity exceeds the Milky Way background field ($g_{\mathrm{gal}} \approx 2 \times 10^{-10}$~m/s$^2$ at $R_\odot = 8$~kpc) only out to $r \approx \sqrt{GM_\odot / g_{\mathrm{gal}}} \approx 9\,800$~AU $\approx 0.05$~pc.
Beyond this radius, the galactic field dominates.

When the kernel is evaluated at any Solar System radius, the vast majority of its weight falls on galactic-scale distances where $g_{\mathrm{bar}} \approx g_{\mathrm{gal}}$.
The nonlocal smoothed field at the Sun's position is therefore set by the galactic baryonic field, not by the Sun's local field:
\begin{equation}
\tilde{g}\big|_{\mathrm{Solar\;System}} \approx g_{\mathrm{gal}} \approx \frac{v_{\mathrm{circ}}^2}{R_\odot} \approx 2.0 \times 10^{-10} \;\text{m/s}^2.
\end{equation}

The resulting extra acceleration is:
\begin{equation}
g_{\mathrm{extra}} = \sqrt{a_0 \cdot g_{\mathrm{gal}}} = \sqrt{6.12 \times 10^{-11} \times 2.0 \times 10^{-10}} \approx 1.1 \times 10^{-10} \;\text{m/s}^2.
\end{equation}

Crucially, $g_{\mathrm{extra}}$ is \textit{constant} across the entire Solar System: it varies only on kiloparsec scales (the scale of the galactic field gradient), not on AU scales.
Explicit numerical evaluation on a Milky Way rotation curve model (Appendix~\ref{app:solar-numerical}, Table~\ref{tab:solar-numerical}) confirms that $\tilde{g}$ is constant to a relative precision of $3.2 \times 10^{-8}$ across $\pm 40$~AU, and the variation in $g_{\mathrm{extra}}$ is $1.8 \times 10^{-18}$~m/s$^2$.

\subsection{Mechanism 2: A constant acceleration is absorbed by free fall}

The Solar System is in free fall in the galactic gravitational field.
By the equivalence principle (Proposition~14.1 of~\cite{Wells2026c}), a uniform external acceleration does not affect the internal dynamics of a freely falling system.
The constant $g_{\mathrm{extra}} \approx 1.1 \times 10^{-10}$~m/s$^2$ affects the Sun's orbit around the galactic center but does \textit{not} produce anomalous precession, semi-major axis drift, or eccentricity evolution in planetary orbits.

The residual effect comes from the galactic tidal gradient of $g_{\mathrm{extra}}$ across the Solar System.
The galactic field gradient is $dg_{\mathrm{gal}}/dR \approx g_{\mathrm{gal}}/R_\odot \approx 8 \times 10^{-31}$~m/s$^2$/m.
The tidal variation of $g_{\mathrm{extra}}$ across the Solar System ($\sim\!40$~AU) is:
\begin{equation}
\Delta g_{\mathrm{extra}} \approx \frac{1}{2}\sqrt{\frac{a_0}{g_{\mathrm{gal}}}} \cdot \frac{g_{\mathrm{gal}}}{R_\odot} \cdot 40\;\mathrm{AU} \approx 1.3 \times 10^{-18} \;\text{m/s}^2.
\end{equation}

This is $\sim\!2 \times 10^{-16}$ of Solar gravity at 1~AU---many orders of magnitude below any measurement threshold.

\subsection{Summary}

\begin{center}
\begin{tabular}{lll}
\toprule
\textbf{Effect} & \textbf{Magnitude} & \textbf{Detectable?} \\
\midrule
$g_{\mathrm{extra}}$ at Sun's position & $1.1 \times 10^{-10}$ m/s$^2$ & No (constant, free fall) \\
Tidal variation across Solar System & $1.3 \times 10^{-18}$ m/s$^2$ & No ($10^{-16}$ of $g_\odot$) \\
Anomalous planetary precession & $\sim\!0$ & No \\
\bottomrule
\end{tabular}
\end{center}

The PPN parameters $\beta = \gamma = 1$~\cite{Wells2026c} are properties of the RST effective metric and remain secure.
The nonlocal correction is a separate, additive term that does not modify the metric structure.
The Solar System tests reported in~\cite{Wells2026c} validate the coordinate-time metric; the nonlocal term derived in this paper is shown here to be independently safe by the mechanisms above.

\textit{Status:~\textbf{S}.}
The safety follows from two arguments (nonlocal averaging at galactic scales; equivalence principle in the freely falling Solar System frame), each supported by explicit numerical computation (Appendix~\ref{app:solar-numerical}).
The residual galactic tidal contribution is ${\sim}10^{-16}$ of Solar gravity.
The status is~\textbf{S} rather than~\textbf{T} because the computation uses a simplified Milky Way rotation curve model rather than a fully realistic mass distribution; the conclusion is expected to be robust to the Milky Way model but has not been verified against detailed galactic mass models.

%===================================================================
\section{Status Discipline}
\label{sec:status-table}
%===================================================================

\begin{table}[h]
\centering
\caption{Derivation status of all model ingredients and empirical claims.}
\label{tab:status}
\begin{tabular}{llp{6.5cm}}
\toprule
\textbf{Claim} & \textbf{Status} & \textbf{Note} \\
\midrule
\multicolumn{3}{l}{\textit{Theorem-derived}} \\
Covering degree $d = 2$ & \textbf{T} & \cite{Wells2026a}, Theorem~7.4 \\
Solar System safety & \textbf{S} & Nonlocal term constant across Solar System (Appendix~\ref{app:solar-numerical}); tidal residual $\sim\!10^{-16}$ of $g_\odot$, \S\ref{sec:solar} \\
\addlinespace
\multicolumn{3}{l}{\textit{Structurally inferred}} \\
$a_0 = a_{\mathrm{nat}} \times N_m / R^{16}$ & \textbf{S} & Dimensional analysis, \S\ref{sec:accel-scale} \\
$L_{\mathrm{GL}} = s_{\mathrm{nat}} \times R^{12}$ & \textbf{S} & Dimensional analysis, \S\ref{sec:coherence} \\
Sqrt-exponential kernel ($p = 0.5$) & \textbf{S} & From $r = \rho^2$, \S\ref{sec:kernel} \\
Gravitational limit existence & \textbf{S} & From P5 + gravitation, \S\ref{sec:grav-limit} \\
$N_m$ as gravitational sector factor & \textbf{S} & Sector assignment, \S\ref{sec:Nm} \\
$R^{16}$ exponent & \textbf{S} & Imported from~\cite{Wells2026c} \\
\addlinespace
\multicolumn{3}{l}{\textit{Implementation-fixed}} \\
Exponential kernel form in $\rho$-space & \textbf{S\!f} & Assumed, \S\ref{sec:kernel} \\
$\Upsilon_{\mathrm{disk}} = 0.5$, $\Upsilon_{\mathrm{bulge}} = 0.7$ & \textbf{S\!f} & SPARC standard \\
Curvature modulation (optional) & \textbf{S\!f} & Non-critical, \S\ref{sec:curvature} \\
\addlinespace
\multicolumn{3}{l}{\textit{Empirically tested}} \\
Median $\chi^2 = 10.71$ competitive with MOND & \textbf{E} & 0 params, \S\ref{sec:global} \\
$d = 2$ sharply confirmed & \textbf{E} & \S\ref{sec:d-scan} \\
$p = 0.5$ confirmed as optimum & \textbf{E} & \S\ref{sec:kernel-scan} \\
Velocity-dependent pattern & \textbf{E} & \S\ref{sec:velocity} \\
Performance ceiling traced to baryonic input & \textbf{E} & \S\ref{sec:ceiling} \\
\addlinespace
\multicolumn{3}{l}{\textit{Open}} \\
$G$ from natural units & \textbf{O} & Required to close hierarchy \\
Mass-dependent gravitational limit & \textbf{O} & Would improve massive spirals \\
Cosmological implications & \textbf{O} & CMB, BAO, lensing \\
Extension to ellipticals and clusters & \textbf{O} & \\
\bottomrule
\end{tabular}
\end{table}

%===================================================================
\section{Conclusion}
\label{sec:conclusion}
%===================================================================

This paper does not claim that the SPARC benchmark proves the compact-fiber framework.
It claims that the framework, as developed in~\cite{Wells2026a,Wells2026b,Wells2026c}, yields a specific galactic rotation model with no calibrated galactic parameters, and that this model performs competitively with standard MOND on the SPARC sample.

The derivation chain:
\begin{quote}
\noindent
Postulates (P1--P5) \\
$\;\to$ Compact fiber~\cite{Wells2026a}: $d = 2$, $N_m = 4$, $N_e = 8$ \\
$\;\to$ Second-power relation~\cite{Wells2026b}: $r = \rho^2$ \\
$\;\to$ Structural constant~\cite{Wells2026b}: $6 = \mathrm{period}_{\mathrm{ref}} + e_{\mathrm{half}}$ \\
$\;\to$ Coordinate-time mechanism~\cite{Wells2026c}: $3v^2/c^2$ \\
$\;\to$ $R^{16}$ hierarchy exponent~\cite{Wells2026c} \\
$\;\to$ Gravitational limit: progression $=$ gravitation at $r_{\mathrm{GL}}$ \\
$\;\to$ Derived scales: $a_0 = a_{\mathrm{nat}} \times N_m / R^{16}$, $L_{\mathrm{GL}} = s_{\mathrm{nat}} \times R^{12}$ \\
$\;\to$ Kernel shape: $\exp(-|\sqrt{r} - \sqrt{r'}|/L)$ from $r = \rho^2$ \\
$\;\to$ Coupling law: $g + \sqrt{a_0 \cdot \tilde{g}}$ from $d = 2$ \\
$\;\to$ 171-galaxy test: median $\chi^2 = 10.71$ (vs.\ MOND $11.25$)
\end{quote}

Every ingredient in the headline model is either imported from~\cite{Wells2026a,Wells2026b,Wells2026c}, structurally inferred from them under stated assumptions, or fixed by the explicitly declared auxiliary assumption and implementation choices recorded in the status table (\S\ref{sec:status-table}).
No parameters are calibrated.
The model's performance signature---dominant in gas-rich dwarfs, competitive in giant spirals, deficit in disk-dominated massive spirals---identifies the mass-dependent gravitational limit as the next structural target.

The model's residuals are dominated by baryonic input uncertainty (galaxy-to-galaxy M/L variation), not by the gravitational theory.
In the mid-to-outer disk, where baryonic-input ambiguities are plausibly reduced relative to the inner disk, RST achieves near-zero mean fractional residuals in both point-pooled and galaxy-level analyses (\S\ref{sec:radial}, Appendix~\ref{app:robustness}), while MOND systematically overpredicts by $3$--$6\%$.

%===================================================================
\appendix
%===================================================================

\section{Reproducibility Protocol}
\label{app:reproducibility}

This appendix specifies every operational detail needed to reproduce the headline results from the SPARC data files and the stated model parameters.

\subsection{Data source and galaxy selection}

The rotation-curve data are from the SPARC Rotmod\_LTG catalog~\cite{Lelli2016}, distributed as individual \texttt{.dat} files (one per galaxy) in the \texttt{Rotmod\_LTG/} directory.
Each file contains 8 whitespace-separated columns:
\begin{center}
\texttt{Rad~~Vobs~~errV~~Vgas~~Vdisk~~Vbul~~SBdisk~~SBbul}
\end{center}
with units kpc, km/s, km/s, km/s, km/s, km/s, $L_\odot$/pc$^2$, $L_\odot$/pc$^2$.

Galaxy inclusion rule: a galaxy is included if, after dropping rows with missing values and rows with $\texttt{errV} \leq 0$, at least 5 data points remain.
This yields 171 galaxies.
No galaxies are excluded by morphological type, data quality flag, or model performance.
Files are sorted alphabetically; the galaxy list is fully determined by the inclusion rule applied to the SPARC distribution.

\subsection{Baryonic velocity construction}

The baryonic circular velocity at each radius is:
\begin{equation}
\label{eq:vbar}
v_{\mathrm{bar}}(r) = \sqrt{\Upsilon_{\mathrm{disk}}\,v_{\mathrm{disk}}^2(r) + \Upsilon_{\mathrm{bulge}}\,v_{\mathrm{bulge}}^2(r) + v_{\mathrm{gas}}^2(r)},
\end{equation}
where $v_{\mathrm{gas}}$, $v_{\mathrm{disk}}$, $v_{\mathrm{bulge}}$ are the SPARC columns \texttt{Vgas}, \texttt{Vdisk}, \texttt{Vbul} respectively, and the mass-to-light ratios are $\Upsilon_{\mathrm{disk}} = 0.5$, $\Upsilon_{\mathrm{bulge}} = 0.7$ (fixed, SPARC standard at 3.6~$\mu$m).
Negative SPARC velocity components (which can arise from counter-rotating gas) are clipped to zero before squaring.
Bulgeless galaxies have $v_{\mathrm{bulge}} = 0$ at all radii.

The baryonic acceleration in SI units is:
\begin{equation}
\label{eq:gbar-SI}
g_{\mathrm{bar}}(r) = \frac{(v_{\mathrm{bar}} \times 10^3)^2}{r \times 3.0857 \times 10^{19}} \quad [\mathrm{m/s}^2],
\end{equation}
where $v_{\mathrm{bar}}$ is in km/s and $r$ is in kpc.

\subsection{RST model evaluation}
\label{app:rst-algorithm}

The RST model is evaluated on the discrete observed radii only (no interpolation).
The algorithm for each galaxy:

\begin{enumerate}
\item Read $\{r_i, v_{\mathrm{obs},i}, \sigma_{v,i}, v_{\mathrm{gas},i}, v_{\mathrm{disk},i}, v_{\mathrm{bulge},i}\}$ for $i = 1, \ldots, N$.
\item Compute $v_{\mathrm{bar},i}$ via Eq.~\eqref{eq:vbar}.
\item Compute $g_{\mathrm{bar},i}$ via Eq.~\eqref{eq:gbar-SI}.
\item Discard any point with non-finite $r$, $g_{\mathrm{bar}}$, or $g_{\mathrm{bar}} \leq 0$. Sort remaining points by ascending~$r$.
\item For each point~$i$, compute:
\begin{equation}
\tilde{g}_i = \frac{\sum_{j} g_{\mathrm{bar},j} \cdot \exp\!\bigl(-|\sqrt{r_i} - \sqrt{r_j}|/L_{\mathrm{GL}}\bigr)}{\sum_{j} \exp\!\bigl(-|\sqrt{r_i} - \sqrt{r_j}|/L_{\mathrm{GL}}\bigr)},
\end{equation}
where the sums run over all $N$ valid points (including $j = i$), and $r$ is in kpc.
\item Compute $g_{\mathrm{total},i} = g_{\mathrm{bar},i} + \sqrt{\max(a_0 \cdot \max(\tilde{g}_i, 0),\; 0)}$.
\item Compute $v_{\mathrm{pred},i} = \sqrt{\max(g_{\mathrm{total},i} \cdot r_i \cdot 3.0857 \times 10^{19},\; 0)} \;/\; 10^3$ \quad [km/s].
\end{enumerate}

Parameters: $a_0 = 6.1225 \times 10^{-11}$~m/s$^2$, $L_{\mathrm{GL}} = 0.31757$~kpc.

\subsection{MOND benchmark evaluation}

The MOND benchmark uses the exponential-style RAR interpolation:
\begin{equation}
\label{eq:mond-alg}
g_{\mathrm{total}} = \frac{g_{\mathrm{bar}}}{1 - \exp\!\bigl(-\!\sqrt{g_{\mathrm{bar}} / a_0^{\mathrm{MOND}}}\bigr)},
\end{equation}
with $a_0^{\mathrm{MOND}} = 1.2 \times 10^{-10}$~m/s$^2$.
The baryonic acceleration $g_{\mathrm{bar}}$ is constructed identically to the RST model (Eq.~\ref{eq:gbar-SI}, same $\Upsilon$ values).
The predicted velocity is $v_{\mathrm{pred}} = \sqrt{g_{\mathrm{total}} \cdot r}$ in the same units.

\subsection{DFT-B ablation comparator}

DFT-B uses the same coupling law as RST ($g_{\mathrm{total}} = g_{\mathrm{bar}} + \sqrt{a_c \cdot \tilde{g}}$) but with a compact triangular kernel:
\begin{equation}
K_{\mathrm{DFT\text{-}B}}(r, r') = \max\!\bigl(0,\; 1 - |r - r'|/L_c\bigr),
\end{equation}
and two calibrated parameters $(a_c, L_c)$ obtained by minimizing the median $\chi^2$/pt on the full 171-galaxy sample via grid search: $a_c \in \{2 \times 10^{-11}, \ldots, 8 \times 10^{-10}\}$ (25 log-spaced values), $L_c \in \{0.3, \ldots, 15.0\}$~kpc (25 linearly spaced values).
The baryonic input construction is identical to the RST model.

\subsection{Metric definitions}

\begin{align}
\chi^2/\mathrm{pt} &= \frac{1}{N}\sum_{i=1}^{N}\left(\frac{v_{\mathrm{obs},i} - v_{\mathrm{pred},i}}{\sigma_{v,i}}\right)^{\!2}, \label{eq:chi2} \\[4pt]
\text{Win (RST vs.\ MOND)} &: \quad \chi^2_{\mathrm{RST}} < \chi^2_{\mathrm{MOND}} \;\text{for that galaxy.} \label{eq:win}
\end{align}

The sign test $p$-value is the two-sided binomial probability of observing at least as extreme a split under the null hypothesis of equal performance.
Ties ($\chi^2_{\mathrm{RST}} = \chi^2_{\mathrm{MOND}}$ to machine precision) do not occur in practice.

\subsection{Scan protocols}

All scans vary one parameter while holding all others at the canonical values ($a_0 = 6.1225 \times 10^{-11}$, $L_{\mathrm{GL}} = 0.31757$~kpc, $p = 0.5$, $d = 2$, $\eta = 0$).

Covering-degree scan: $d \in \{1.5, 1.8, 2.0, 2.2, 2.5, 3.0\}$.

Kernel-power scan: $p \in \{0.30, 0.40, 0.50, 0.60, 0.70, 0.80, 1.00\}$.

Curvature-modulation scan: $\eta \in \{0.00, 0.10, 0.20, 0.30, 0.45, 0.55, 0.70\}$.

Mass-to-light scan: $\Upsilon_{\mathrm{disk}} \in \{0.3, 0.4, 0.5, 0.6, 0.7\}$ with $\Upsilon_{\mathrm{bulge}} = 0.7$ fixed.

\subsection{Radial residual binning}

The fractional residual at each data point is $(v_{\mathrm{pred}} - v_{\mathrm{obs}})/v_{\mathrm{obs}}$.
For the radial profile (Table~\ref{tab:radial}), each data point's normalized radius is $r/r_{\max}$ where $r_{\max}$ is the largest observed radius in that galaxy.
Points are pooled across all 171 galaxies and binned into five equal-width bins: $[0, 0.2)$, $[0.2, 0.4)$, $[0.4, 0.6)$, $[0.6, 0.8)$, $[0.8, 1.0]$.
Mean and rms are computed within each bin.

\subsection{$V_{\mathrm{flat}}$ definition}

For the velocity-bin analysis (Table~\ref{tab:velocity}), $V_{\mathrm{flat}}$ is defined as the median of $v_{\mathrm{obs}}$ over the outer quarter of each galaxy's radial range (the last $\max(3, N/4)$ data points, sorted by radius).

\subsection{Code and data availability}

The complete analysis code is available at\\\texttt{github.com/arwells-research/compact-fiber-galactic-rotation}.
The repository contains:
the SPARC Rotmod\_LTG data files under \texttt{data/sparc/Rotmod\_LTG/} with SHA-256 checksum verification;
the RST model implementation in \texttt{src/rst\_model.py};
the MOND and DFT-B implementations in \texttt{src/models.py};
a machine-readable CSV file of per-galaxy results (Appendix~\ref{app:galaxies}) as \texttt{rst\_sparc\_results.csv};
and experiment scripts in \texttt{experiments/} that regenerate all tables in this paper.

To reproduce the headline results from the raw SPARC data:
\begin{verbatim}
  cd experiments && python run_rst_final.py
\end{verbatim}
This single command reads the SPARC files, applies the inclusion rule (\S\ref{app:reproducibility}), evaluates RST, MOND, and DFT-B on all qualifying galaxies, and outputs the summary statistics, scan tables, and per-galaxy CSV.

\section{Galaxy-Level Residual Robustness Check}
\label{app:robustness}

The radial residual profile in Table~\ref{tab:radial} pools all per-point measurements across 171 galaxies, which could be dominated by larger or more densely sampled galaxies.
As a robustness check, this appendix computes the same radial profile using galaxy-level averaging: for each galaxy, the mean fractional residual is computed within each radial bin, and the median across galaxies is reported.

\begin{table}[h]
\centering
\caption{Galaxy-level median of per-galaxy mean fractional residuals, by normalized radius. Compare with the point-pooled Table~\ref{tab:radial}.}
\label{tab:radial-robust}
\begin{tabular}{rrrr}
\toprule
$r/r_{\max}$ & RST median & MOND median & $n_{\mathrm{gal}}$ \\
\midrule
$0.0$--$0.2$ & $+8.1\%$ & $+7.4\%$ & 124 \\
$0.2$--$0.4$ & $+0.1\%$ & $+0.3\%$ & 139 \\
$0.4$--$0.6$ & $-2.5\%$ & $+1.5\%$ & 136 \\
$0.6$--$0.8$ & $-2.7\%$ & $+3.0\%$ & 126 \\
$0.8$--$1.0$ & $+1.8\%$ & $+6.1\%$ & 160 \\
\bottomrule
\end{tabular}
\end{table}

The galaxy-level analysis confirms the point-pooled pattern: RST achieves near-zero or slightly negative median residuals in the mid-to-outer disk ($0.2$--$0.8\,r_{\max}$: medians $-2.7\%$ to $+0.1\%$), while MOND systematically overpredicts ($+0.3\%$ to $+6.1\%$).
The conclusion is robust to the averaging method.

\section{DFT-B Calibrated Parameters}
\label{app:dftb}

The DFT-B ablation comparator was calibrated by grid search over $a_c \in \{2.0 \times 10^{-11}, \ldots, 8.0 \times 10^{-10}\}$ (25 log-spaced values) and $L_c \in \{0.3, \ldots, 15.0\}$~kpc (25 linearly spaced values), minimizing the median $\chi^2$/pt across all 171 galaxies.

Selected best-fit parameters: $a_c = 5.87 \times 10^{-11}$~m/s$^2$, $L_c = 11.32$~kpc.
Resulting median $\chi^2$/pt $= 10.13$.

\section{Solar System Numerical Computation}
\label{app:solar-numerical}

This appendix provides the explicit numerical evaluation supporting the Solar System safety argument of~\S\ref{sec:solar}.

The computation evaluates $\tilde{g}$ at the Sun's galactic position ($R = 8$~kpc) using a flat Milky Way rotation curve ($v_{\mathrm{circ}} = 220$~km/s for $R > 1$~kpc, rising linearly for $R < 1$~kpc) discretized on 600 radial points from $R = 0.01$ to $R = 30$~kpc.
The RST kernel $K(R, R') = \exp(-|\sqrt{R} - \sqrt{R'}|/L_{\mathrm{GL}})$ is evaluated with $L_{\mathrm{GL}} = 0.31757$~kpc.

\begin{table}[h]
\centering
\caption{$\tilde{g}$ and $g_{\mathrm{extra}}$ at six positions spanning the Solar System, evaluated on the galactic rotation curve.
All positions are within $\pm 40$~AU of $R_\odot = 8$~kpc.}
\label{tab:solar-numerical}
\begin{tabular}{rrr}
\toprule
\textbf{Offset from $R_\odot$} & $\tilde{g}$ (m/s$^2$) & $g_{\mathrm{extra}}$ (m/s$^2$) \\
\midrule
$-10$ AU & $2.0294 \times 10^{-10}$ & $1.1147 \times 10^{-10}$ \\
$-1$ AU & $2.0294 \times 10^{-10}$ & $1.1147 \times 10^{-10}$ \\
$0$ (exact) & $2.0294 \times 10^{-10}$ & $1.1147 \times 10^{-10}$ \\
$+1$ AU & $2.0294 \times 10^{-10}$ & $1.1147 \times 10^{-10}$ \\
$+10$ AU & $2.0294 \times 10^{-10}$ & $1.1147 \times 10^{-10}$ \\
$+40$ AU & $2.0294 \times 10^{-10}$ & $1.1147 \times 10^{-10}$ \\
\bottomrule
\end{tabular}
\end{table}

The maximum relative variation in $\tilde{g}$ across $\pm 40$~AU is $3.2 \times 10^{-8}$.
The variation in $g_{\mathrm{extra}}$ is $1.8 \times 10^{-18}$~m/s$^2$, which is $2.4 \times 10^{-16}$ of Solar gravity at 1~AU ($5.93 \times 10^{-3}$~m/s$^2$).

The kernel weight decomposition at $R = 8$~kpc:
\begin{center}
\begin{tabular}{lrr}
\toprule
\textbf{Region} & \textbf{Weight fraction} & \textbf{Contribution to $\tilde{g}$} \\
\midrule
$R < 1$ kpc (inner galaxy) & 0.22\% & 1.2\% \\
$1$--$15$ kpc (near Sun) & 97.0\% & 97.6\% \\
$R > 15$ kpc (outer galaxy) & 2.8\% & 1.2\% \\
\bottomrule
\end{tabular}
\end{center}

The dominant contribution to $\tilde{g}$ comes from the galactic field in the $1$--$15$~kpc range around the Sun's position; the inner galaxy and outer galaxy each contribute ${\sim}1\%$.
The Sun's own gravitational field does not appear in this computation because the galactic rotation curve describes the smooth baryonic field of the disk, not individual stellar fields.
The RST model operates on the galactic-scale baryonic profile; the Sun is one of ${\sim}10^{11}$ stars contributing to that smooth field.

The galactic tidal gradient of the extra acceleration:
\begin{equation}
\frac{dg_{\mathrm{extra}}}{dR} = \frac{1}{2}\sqrt{\frac{a_0}{\tilde{g}}} \cdot \frac{d\tilde{g}}{dR} \approx 2.4 \times 10^{-31} \;\text{m/s}^2\text{/m}.
\end{equation}
Over 40~AU ($6 \times 10^{12}$~m): $\Delta g_{\mathrm{extra}} \approx 1.4 \times 10^{-18}$~m/s$^2$.

This computation can be reproduced using the script \texttt{experiments/run\_solar\_safety.py} in the repository.

\section{Per-Galaxy Results}
\label{app:galaxies}

\begin{longtable}{lrrrrrl}
\caption{Per-galaxy $\chi^2$/pt for the RST pure structural model and MOND baseline on 171 SPARC galaxies.}
\label{tab:galaxies}\\
\toprule
\textbf{Galaxy} & $N$ & $V_{\mathrm{flat}}$ & $\chi^2_{\mathrm{RST}}$ & $\chi^2_{\mathrm{MOND}}$ & $\log_{10}$ ratio & Winner \\
\midrule
\endfirsthead
\toprule
\textbf{Galaxy} & $N$ & $V_{\mathrm{flat}}$ & $\chi^2_{\mathrm{RST}}$ & $\chi^2_{\mathrm{MOND}}$ & $\log_{10}$ ratio & Winner \\
\midrule
\endhead
\bottomrule
\endfoot
CamB                 &   9 &   16.8 &    79.68 &   108.20 & -0.133 & RST \\
D564-8               &   6 &   24.0 &     7.72 &    18.30 & -0.375 & RST \\
D631-7               &  16 &   58.0 &    16.75 &    30.06 & -0.254 & RST \\
DDO064               &  14 &   46.4 &     0.34 &     0.36 & -0.034 & RST \\
DDO154               &  12 &   47.4 &    40.05 &    89.15 & -0.347 & RST \\
DDO161               &  31 &   66.8 &    67.26 &   149.63 & -0.347 & RST \\
DDO168               &  10 &   54.8 &    16.64 &    30.06 & -0.257 & RST \\
DDO170               &   8 &   60.0 &    45.50 &   152.46 & -0.525 & RST \\
ESO079-G014          &  15 &  177.0 &    15.30 &     7.99 & +0.282 & MOND \\
ESO116-G012          &  15 &  111.0 &    24.06 &    12.89 & +0.271 & MOND \\
ESO444-G084          &   7 &   62.7 &    41.10 &    27.78 & +0.170 & MOND \\
ESO563-G021          &  30 &  315.0 &    46.84 &    53.53 & -0.058 & RST \\
F561-1               &   6 &   50.4 &    27.21 &    40.20 & -0.169 & RST \\
F563-1               &  17 &  111.2 &     4.81 &     2.78 & +0.238 & MOND \\
F563-V1              &   6 &   28.6 &    23.23 &    34.37 & -0.170 & RST \\
F563-V2              &  10 &  118.0 &     5.95 &     4.33 & +0.138 & MOND \\
F565-V2              &   7 &   78.5 &     3.64 &     1.09 & +0.522 & MOND \\
F567-2               &   5 &   48.1 &     2.98 &     7.11 & -0.377 & RST \\
F568-1               &  12 &  133.0 &     9.71 &     6.81 & +0.154 & MOND \\
F568-3               &  18 &  112.0 &     2.68 &     2.78 & -0.016 & RST \\
F568-V1              &  15 &  113.0 &    12.83 &     9.39 & +0.135 & MOND \\
F571-8               &  13 &  140.0 &    43.88 &    40.57 & +0.034 & MOND \\
F571-V1              &   7 &   83.9 &     0.25 &     0.78 & -0.501 & RST \\
F574-1               &  14 &   99.5 &     8.18 &     3.52 & +0.367 & MOND \\
F574-2               &   5 &   36.0 &    15.82 &    23.29 & -0.168 & RST \\
F579-V1              &  14 &  114.0 &     4.50 &     3.42 & +0.119 & MOND \\
F583-1               &  25 &   82.5 &     6.05 &     3.48 & +0.240 & MOND \\
F583-4               &  12 &   67.2 &     0.09 &     1.31 & -1.153 & RST \\
IC2574               &  34 &   64.5 &   102.53 &   238.45 & -0.367 & RST \\
IC4202               &  32 &  244.0 &    44.97 &    37.41 & +0.080 & MOND \\
KK98-251             &  15 &   34.2 &    15.65 &    29.41 & -0.274 & RST \\
NGC0024              &  29 &  107.0 &    11.81 &    10.46 & +0.053 & MOND \\
NGC0055              &  21 &   86.2 &    25.36 &    51.14 & -0.305 & RST \\
NGC0100              &  21 &   89.0 &     1.50 &     1.99 & -0.122 & RST \\
NGC0247              &  26 &  105.5 &    13.02 &     4.36 & +0.475 & MOND \\
NGC0289              &  28 &  166.0 &     3.48 &    11.62 & -0.524 & RST \\
NGC0300              &  25 &   94.2 &     7.37 &     2.40 & +0.487 & MOND \\
NGC0801              &  13 &  216.0 &    59.56 &    86.90 & -0.164 & RST \\
NGC0891              &  18 &  213.0 &    91.15 &    34.81 & +0.418 & MOND \\
NGC1003              &  36 &  112.0 &     6.86 &    13.08 & -0.281 & RST \\
NGC1090              &  24 &  162.0 &    22.23 &    44.28 & -0.299 & RST \\
NGC1705              &  14 &   71.5 &    15.78 &    14.54 & +0.035 & MOND \\
NGC2366              &  26 &   49.7 &    10.15 &    25.37 & -0.398 & RST \\
NGC2403              &  73 &  134.0 &   106.79 &    75.59 & +0.150 & MOND \\
NGC2683              &  11 &  155.0 &     1.09 &     2.14 & -0.292 & RST \\
NGC2841              &  50 &  282.5 &    98.18 &    86.26 & +0.056 & MOND \\
NGC2903              &  34 &  181.5 &    50.26 &    12.60 & +0.601 & MOND \\
NGC2915              &  30 &   84.7 &    10.43 &     7.61 & +0.137 & MOND \\
NGC2955              &  24 &  255.5 &    24.13 &     7.96 & +0.482 & MOND \\
NGC2976              &  27 &   83.8 &     3.42 &     3.04 & +0.051 & MOND \\
NGC2998              &  13 &  213.0 &     3.84 &    15.47 & -0.605 & RST \\
NGC3109              &  25 &   66.2 &    11.33 &     5.37 & +0.324 & MOND \\
NGC3198              &  43 &  148.5 &     1.65 &    13.42 & -0.911 & RST \\
NGC3521              &  41 &  210.0 &     1.78 &     0.73 & +0.387 & MOND \\
NGC3726              &  12 &  168.0 &    24.86 &    24.98 & -0.002 & RST \\
NGC3741              &  21 &   50.1 &     4.86 &     1.35 & +0.556 & MOND \\
NGC3769              &  12 &  118.0 &     4.13 &     4.66 & -0.053 & RST \\
NGC3877              &  13 &  170.0 &    16.82 &    13.12 & +0.108 & MOND \\
NGC3893              &  10 &  176.0 &     1.58 &     0.89 & +0.247 & MOND \\
NGC3917              &  17 &  137.0 &     4.03 &     2.86 & +0.150 & MOND \\
NGC3949              &   7 &  165.0 &     6.93 &     2.26 & +0.486 & MOND \\
NGC3953              &   8 &  223.0 &     2.70 &     2.48 & +0.038 & MOND \\
NGC3972              &  10 &  134.0 &     2.70 &     1.90 & +0.153 & MOND \\
NGC3992              &   9 &  241.0 &     9.88 &     9.24 & +0.029 & MOND \\
NGC4010              &  12 &  124.0 &     3.33 &     3.99 & -0.078 & RST \\
NGC4013              &  36 &  178.0 &    10.13 &     7.49 & +0.131 & MOND \\
NGC4051              &   7 &  154.0 &    14.47 &    11.42 & +0.103 & MOND \\
NGC4068              &   6 &   36.0 &    13.69 &    24.49 & -0.253 & RST \\
NGC4085              &   7 &  133.0 &    18.96 &    13.13 & +0.160 & MOND \\
NGC4088              &  12 &  171.0 &    45.84 &    35.89 & +0.106 & MOND \\
NGC4100              &  24 &  158.5 &     2.73 &     3.39 & -0.094 & RST \\
NGC4138              &   7 &  147.0 &     0.72 &     1.21 & -0.229 & RST \\
NGC4157              &  17 &  185.5 &    11.66 &     9.09 & +0.108 & MOND \\
NGC4183              &  23 &  110.0 &     1.53 &     3.69 & -0.382 & RST \\
NGC4214              &  14 &   80.6 &     5.45 &     5.67 & -0.017 & RST \\
NGC4217              &  19 &  177.5 &    86.38 &    55.71 & +0.190 & MOND \\
NGC4389              &   6 &   95.9 &    62.99 &    54.01 & +0.067 & MOND \\
NGC4559              &  32 &  121.5 &     7.52 &    13.89 & -0.266 & RST \\
NGC5005              &  18 &  264.5 &     3.26 &     0.54 & +0.778 & MOND \\
NGC5033              &  22 &  193.0 &    19.67 &    13.86 & +0.152 & MOND \\
NGC5055              &  28 &  179.0 &  1365.96 &   770.51 & +0.249 & MOND \\
NGC5371              &  19 &  208.5 &   172.14 &   170.72 & +0.004 & MOND \\
NGC5585              &  24 &   90.8 &    22.01 &    26.21 & -0.076 & RST \\
NGC5907              &  19 &  215.5 &     9.90 &    33.57 & -0.530 & RST \\
NGC5985              &  33 &  292.5 &   154.96 &   158.40 & -0.010 & RST \\
NGC6015              &  44 &  155.0 &    27.65 &    32.01 & -0.064 & RST \\
NGC6195              &  23 &  254.0 &    54.13 &    26.57 & +0.309 & MOND \\
NGC6503              &  31 &  115.0 &    17.80 &     2.81 & +0.802 & MOND \\
NGC6674              &  15 &  242.0 &    74.80 &    24.53 & +0.484 & MOND \\
NGC6946              &  58 &  158.0 &    55.37 &    25.66 & +0.334 & MOND \\
NGC7331              &  36 &  239.0 &    65.88 &    36.00 & +0.262 & MOND \\
NGC7793              &  46 &  109.0 &     1.27 &     0.91 & +0.145 & MOND \\
NGC7814              &  18 &  214.0 &    18.60 &    27.26 & -0.166 & RST \\
PGC51017             &   6 &   18.3 &   105.22 &   131.44 & -0.097 & RST \\
UGC00128             &  22 &  132.0 &   212.43 &    96.97 & +0.341 & MOND \\
UGC00191             &   9 &   81.0 &     8.85 &    29.68 & -0.526 & RST \\
UGC00731             &  12 &   73.9 &    18.15 &    16.36 & +0.045 & MOND \\
UGC00891             &   5 &   60.1 &     7.18 &    47.48 & -0.820 & RST \\
UGC01230             &  11 &  103.0 &     2.68 &     3.69 & -0.138 & RST \\
UGC01281             &  25 &   55.4 &     0.74 &     1.12 & -0.180 & RST \\
UGC02023             &   5 &   47.9 &     1.56 &     2.56 & -0.215 & RST \\
UGC02259             &   8 &   88.3 &    85.57 &    45.61 & +0.273 & MOND \\
UGC02455             &   8 &   53.1 &   126.54 &   125.21 & +0.005 & MOND \\
UGC02487             &  17 &  331.0 &   362.92 &   261.40 & +0.143 & MOND \\
UGC02885             &  19 &  298.0 &     3.36 &     2.83 & +0.074 & MOND \\
UGC02916             &  43 &  209.5 &   225.66 &    75.87 & +0.473 & MOND \\
UGC02953             & 115 &  264.0 &   125.96 &   200.40 & -0.202 & RST \\
UGC03205             &  48 &  219.0 &     8.25 &    12.14 & -0.168 & RST \\
UGC03546             &  30 &  193.0 &    48.40 &    12.24 & +0.597 & MOND \\
UGC03580             &  47 &  122.0 &   125.72 &    87.56 & +0.157 & MOND \\
UGC04278             &  25 &   88.2 &     3.98 &     2.78 & +0.155 & MOND \\
UGC04305             &  22 &   35.1 &    50.92 &    68.66 & -0.130 & RST \\
UGC04325             &   8 &   91.6 &    40.79 &    29.35 & +0.143 & MOND \\
UGC04483             &   8 &   24.2 &     5.29 &    11.25 & -0.328 & RST \\
UGC04499             &   9 &   73.3 &     2.47 &     9.57 & -0.588 & RST \\
UGC05005             &  11 &   99.4 &     0.76 &     2.46 & -0.513 & RST \\
UGC05253             &  73 &  243.5 &    45.07 &     6.91 & +0.814 & MOND \\
UGC05414             &   6 &   56.8 &     4.25 &    11.13 & -0.418 & RST \\
UGC05716             &  12 &   73.0 &    89.14 &    37.51 & +0.376 & MOND \\
UGC05721             &  23 &   78.6 &    13.43 &    12.32 & +0.037 & MOND \\
UGC05750             &  11 &   77.6 &     1.15 &     3.21 & -0.445 & RST \\
UGC05764             &  10 &   50.0 &   397.87 &   219.66 & +0.258 & MOND \\
UGC05829             &  11 &   65.6 &     0.70 &     2.85 & -0.609 & RST \\
UGC05918             &   8 &   42.6 &     0.72 &     1.21 & -0.227 & RST \\
UGC05986             &  15 &  109.0 &    44.70 &    30.09 & +0.172 & MOND \\
UGC05999             &   5 &   97.7 &     2.06 &     5.80 & -0.449 & RST \\
UGC06399             &   9 &   85.6 &     3.38 &     0.77 & +0.644 & MOND \\
UGC06446             &  17 &   83.8 &     9.71 &     4.88 & +0.299 & MOND \\
UGC06614             &  13 &  204.0 &    11.16 &     9.89 & +0.053 & MOND \\
UGC06628             &   7 &   42.3 &    11.96 &    16.36 & -0.136 & RST \\
UGC06667             &   9 &   85.0 &    43.03 &    26.66 & +0.208 & MOND \\
UGC06786             &  45 &  226.0 &   159.89 &   169.81 & -0.026 & RST \\
UGC06787             &  71 &  247.0 &   139.91 &   200.02 & -0.155 & RST \\
UGC06818             &   8 &   71.2 &     5.47 &     8.37 & -0.185 & RST \\
UGC06917             &  11 &  110.0 &     4.39 &     1.44 & +0.485 & MOND \\
UGC06923             &   6 &   79.3 &     1.38 &     2.21 & -0.203 & RST \\
UGC06930             &  10 &  108.0 &     1.12 &     2.53 & -0.354 & RST \\
UGC06973             &   9 &  178.0 &   152.71 &    90.98 & +0.225 & MOND \\
UGC06983             &  17 &  108.5 &     7.92 &     3.62 & +0.340 & MOND \\
UGC07089             &  12 &   77.6 &     4.55 &    10.02 & -0.342 & RST \\
UGC07125             &  13 &   65.1 &    77.73 &   142.51 & -0.263 & RST \\
UGC07151             &  11 &   74.0 &     2.98 &     3.78 & -0.103 & RST \\
UGC07261             &   7 &   74.7 &     0.83 &     0.66 & +0.095 & MOND \\
UGC07323             &  10 &   82.8 &     2.47 &     5.77 & -0.369 & RST \\
UGC07399             &  10 &  103.0 &    96.67 &    76.40 & +0.102 & MOND \\
UGC07524             &  31 &   80.3 &     2.14 &     5.30 & -0.394 & RST \\
UGC07559             &   7 &   31.1 &     5.49 &    10.49 & -0.281 & RST \\
UGC07577             &   9 &   15.9 &    11.35 &    17.57 & -0.190 & RST \\
UGC07603             &  12 &   63.0 &     9.30 &     5.63 & +0.218 & MOND \\
UGC07608             &   8 &   66.0 &     2.99 &     1.31 & +0.358 & MOND \\
UGC07690             &   7 &   56.9 &     2.35 &     2.05 & +0.059 & MOND \\
UGC07866             &   7 &   30.6 &     1.64 &     3.94 & -0.382 & RST \\
UGC08286             &  17 &   83.7 &    52.01 &    27.17 & +0.282 & MOND \\
UGC08490             &  30 &   77.6 &    11.39 &     5.79 & +0.293 & MOND \\
UGC08550             &  11 &   57.5 &     6.93 &     2.57 & +0.431 & MOND \\
UGC08699             &  41 &  181.5 &     4.28 &     6.44 & -0.178 & RST \\
UGC08837             &   8 &   46.6 &    23.61 &    40.32 & -0.232 & RST \\
UGC09037             &  22 &  152.0 &    43.13 &    45.64 & -0.025 & RST \\
UGC09133             &  68 &  228.0 &    27.25 &    68.34 & -0.399 & RST \\
UGC09992             &   5 &   33.4 &     4.35 &     7.32 & -0.225 & RST \\
UGC10310             &   7 &   72.4 &     1.36 &     2.48 & -0.261 & RST \\
UGC11455             &  36 &  272.0 &    11.53 &     6.79 & +0.230 & MOND \\
UGC11557             &  12 &   81.6 &    17.46 &    20.69 & -0.074 & RST \\
UGC11820             &  10 &   81.2 &    19.67 &   132.95 & -0.830 & RST \\
UGC11914             &  65 &  283.5 &    10.01 &    69.81 & -0.843 & RST \\
UGC12506             &  31 &  230.0 &     9.81 &     7.90 & +0.094 & MOND \\
UGC12632             &  15 &   73.1 &     3.81 &     7.10 & -0.271 & RST \\
UGC12732             &  16 &   93.2 &     4.70 &     3.78 & +0.095 & MOND \\
UGCA281              &   7 &   28.8 &     0.80 &     3.35 & -0.625 & RST \\
UGCA442              &   8 &   56.5 &     2.75 &     4.44 & -0.207 & RST \\
UGCA444              &  36 &   35.6 &     0.39 &     2.01 & -0.716 & RST \\
\end{longtable}


%===================================================================
\begin{thebibliography}{99}
%===================================================================

\bibitem{Wells2026a}
A.~R. Wells,
``The compact fiber as the canonical admissibility structure for discrete scalar motion: a necessity result for the Reciprocal System,''
Zenodo, 2026.
\texttt{doi:10.5281/zenodo.19423877}.

\bibitem{Wells2026b}
A.~R. Wells,
``Atomic structure from the compact fiber: transport, screening, and spectroscopic thresholds in the Reciprocal System,''
Zenodo, 2026.
\texttt{doi:10.5281/zenodo.19426034}.

\bibitem{Wells2026c}
A.~R. Wells,
``Effective quantum and gravitational structure from the compact fiber: cross-frame projection, Hilbert representation, Born rule, and coordinate-time metric,''
Zenodo, 2026.
\texttt{doi:10.5281/zenodo.19434652}.

\bibitem{Larson1979}
D.~B. Larson,
\emph{Nothing But Motion},
North Pacific Publishers, Portland, OR, 1979.

\bibitem{Milgrom1983}
M.~Milgrom,
``A modification of the Newtonian dynamics as a possible alternative to the hidden mass hypothesis,''
\emph{Astrophys. J.} \textbf{270}, 365--370, 1983.

\bibitem{Lelli2016}
F.~Lelli, S.~S. McGaugh, and J.~M. Schombert,
``SPARC: Mass Models for 175 Disk Galaxies with Spitzer Photometry and Accurate Rotation Curves,''
\emph{Astron. J.} \textbf{152}, 157, 2016.

\bibitem{McGaugh2016}
S.~S. McGaugh, F.~Lelli, and J.~M. Schombert,
``Radial Acceleration Relation in Rotationally Supported Galaxies,''
\emph{Phys. Rev. Lett.} \textbf{117}, 201101, 2016.

\bibitem{Nehru1985}
K.~V.~K. Nehru,
``Precession of the Planetary Perihelia Due to the Co-ordinate Time,''
\emph{Reciprocity} \textbf{XIV}(1), 11--13, 1985.

\bibitem{Will2014}
C.~M. Will,
``The Confrontation between General Relativity and Experiment,''
\emph{Living Rev. Relativ.} \textbf{17}, 4, 2014.

\bibitem{Bertotti2003}
B.~Bertotti, L.~Iess, and P.~Tortora,
``A test of general relativity using radio links with the Cassini spacecraft,''
\emph{Nature} \textbf{425}, 374--376, 2003.

\bibitem{LIGO2017}
B.~P. Abbott et al.\ (LIGO/Virgo Collaboration),
``Gravitational Waves and Gamma-Rays from a Binary Neutron Star Merger: GW170817 and GRB 170817A,''
\emph{Astrophys. J. Lett.} \textbf{848}, L13, 2017.

\end{thebibliography}

\end{document}
